% ============================================================================
%  BLOQUE III — DIAPOSITIVAS 8–11: RESTRICCIONES Y METODOLOGÍA
% ============================================================================
\section{Restricciones y Metodología}

% ----------------------------------------------------------------------------
%  DIAPOSITIVA 8 — FACTORES DE DISEÑO
% ----------------------------------------------------------------------------
\begin{frame}{Restricciones y Factores de Diseño}
  \begin{columns}[T]
    \column{0.5\textwidth}
    \begin{alertblock}{Restricción temporal y económica}
      \footnotesize
      \begin{itemize}
        \item Un curso académico $\rightarrow$ MVP con funcionalidades troncales
        \item Autor: estudiante sin ingresos — presupuesto \textbf{56,50 €}
        \item Ingeniería inversa como alternativa a la inversión
      \end{itemize}
    \end{alertblock}

    \column{0.5\textwidth}
    \begin{exampleblock}{Impacto positivo}
      \footnotesize
      \begin{itemize}
        \item Obliga a desarrollar competencias transferibles
        \item GATT Xiaomi sin autenticación: una herramienta cara lo habría ocultado
        \item Domótica inteligente $\neq$ grandes inversiones
      \end{itemize}
    \end{exampleblock}

    \begin{block}{Evolución de infraestructura}
      \footnotesize
      \textbf{Fase 1:} servidores UCO con GPU (prototipado)\\
      \textbf{Fase 2:} torre Ryzen 5, 16 GB RAM, 100\% local
    \end{block}
  \end{columns}
\end{frame}

% ----------------------------------------------------------------------------
%  DIAPOSITIVA 9 — PERÍMETRO AIR-GAPPED
% ----------------------------------------------------------------------------
\begin{frame}{Arquitectura de Red Air-Gapped}
  \begin{columns}[T]
    \column{0.5\textwidth}
    \footnotesize
    \textbf{Problema:} los dispositivos comerciales intentan abrir canales hacia sus nubes.

    \textbf{Solución: Router sin salida WAN}
    \begin{itemize}
      \item TP-Link TL-MR6400 sin tarjeta SIM
      \item Actúa como \textit{switch} Ethernet + AP WLAN
      \item Segmento \texttt{192.168.1.0/24} completamente aislado
      \item IPs estáticas — aislamiento físico verificable
    \end{itemize}

    \column{0.5\textwidth}
    \begin{exampleblock}{Propiedades del segmento aislado}
      \footnotesize
      \begin{itemize}
        \item \textbf{Topología plana:} todos los dispositivos en \texttt{192.168.1.0/24}
        \item \textbf{IPs estáticas} sin DHCP — sin dependencias externas
        \item \textbf{Un solo router:} TL-MR6400 sin tarjeta SIM
        \item \textbf{Sin puerta de enlace WAN:} el tráfico no puede salir
        \item \textbf{Verificable:} retirar la SIM demuestra el diseño
      \end{itemize}
    \end{exampleblock}
  \end{columns}
\end{frame}

% ----------------------------------------------------------------------------
%  DIAPOSITIVA 10 — STACK POLÍGLOTA
% ----------------------------------------------------------------------------
\begin{frame}{Stack Políglota y Justificación}
  \begin{block}{Principio}
    \small
    La naturaleza dual de la domótica con IA (tiempo real + flexibilidad cognitiva)
    hace inviable un sistema monolítico en un solo lenguaje.
  \end{block}

  \begin{columns}[T]
    \column{0.33\textwidth}
    \begin{exampleblock}{Python 3.12}
      \footnotesize
      \textbf{Drivers de dispositivos}\\[0.05cm]
      \texttt{tinytuya}, \texttt{bleak}\\[0.1cm]
      \textbf{Enrutador de modelos}\\[0.05cm]
      HTTP server, comfort advisors\\[0.1cm]
      $\rightarrow$ Ecosistema maduro, prototipado rápido
    \end{exampleblock}

    \column{0.33\textwidth}
    \begin{exampleblock}{Node.js 24}
      \footnotesize
      \textbf{Orquestador OpenClaw} \smallgray{(externo)}\\[0.05cm]
      Canales Telegram y dashboard web\\[0.1cm]
      Interfaz de agentes IA\\[0.05cm]
      \texttt{exec} $\rightarrow$ drivers Python\\[0.1cm]
      $\rightarrow$ Gestión nativa de canales
    \end{exampleblock}

    \column{0.33\textwidth}
    \begin{block}{C++20}
      \footnotesize
      \textbf{Firmware ESP32}\\[0.05cm]
      Control determinista de memoria\\[0.1cm]
      Bus LIN \smallgray{(pendiente)}\\[0.1cm]
      $\rightarrow$ \textit{Hard real-time} cuando se necesita
    \end{block}
  \end{columns}

  \centering
  \footnotesize
  Cada lenguaje en su capa óptima, sin acoplar tiempos de ejecución ni forzar compromisos.
\end{frame}

% ----------------------------------------------------------------------------
%  DIAPOSITIVA 11 — METODOLOGÍA DOCS-AS-CODE
% ----------------------------------------------------------------------------
\begin{frame}[fragile]{Metodología: Docs-as-Code y Git}
  \begin{columns}[T]
    \column{0.5\textwidth}
    \small
    \textbf{Control de versiones}
    \begin{itemize}
      \item Repositorio privado GitHub
      \item Cada contribución auditada
      \item Trazabilidad: decisión de diseño $\leftrightarrow$ código
      \item Base para futura iniciativa empresarial
    \end{itemize}

    \bigskip
    \textbf{Docs as Code}
    \begin{itemize}
      \item Memoria \LaTeX{} + código en mismo repo
      \item Consistencia tipográfica profesional
      \item 9 ADR (\textit{Architecture Decision Records})
      \item Diagramas TikZ versionados
    \end{itemize}

    \column{0.5\textwidth}
    \footnotesize
    \textbf{Uso de IA en el desarrollo}
    \begin{itemize}
      \item \textbf{OpenCode:} asistente de codificación (software libre)
      \item \textbf{DeepSeek V4:} modelo de inferencia vía API
      \item \alert{Toda línea revisada y validada por el autor}
    \end{itemize}

    \begin{exampleblock}{Herramientas de gestión}
      SSH con llaves RSA, Cockpit (puerto 9090), \texttt{systemd}.
    \end{exampleblock}
  \end{columns}
\end{frame}
